%!TEX root = forallxyyc.tex
\part{Lógica proposicional verofuncional}
\label{ch.TFL}
\addtocontents{toc}{\protect\mbox{}\protect\hrulefill\par}

\chapter{Primeiros passos para a simbolização}

\section{Validade em virtude da forma}\label{s:ValidityInVirtueOfForm}
Considere este argumento:
	\begin{earg}
		\item[] Está chovendo lá fora.
		\item[] Se está chovendo lá fora, então Jenny está infeliz.
		\item[\texttherefore] Jenny está infeliz.
	\end{earg}
e outro argumento:
	\begin{earg}
		\item[] Jenny é anarco-sindicalista.
		\item[] Se Jenny é anarco-sindicalista, então Dipan é leitor ávido de Tolstói.
		\item[\texttherefore] Dipan é leitor ávido de Tolstói.
	\end{earg}
Ambos os argumentos são válidos, e há um sentido direto em que podemos dizer que eles compartilham uma estrutura comum. Podemos expressar a estrutura da seguinte maneira:
	\begin{earg}
		\item[] $A$
		\item[] Se $A$, então $C$
		\item[\texttherefore] $C$
	\end{earg}
Esta parece uma excelente \emph{estrutura} de argumento. Com efeito, certamente todo argumento com essa estrutura será válido. E essa não é a única boa estrutura de argumento. Considere um argumento como:
	\begin{earg}
		\item[] Jenny está feliz ou triste.
		\item[] Jenny não está feliz.
		\item[\texttherefore] Jenny está triste.
	\end{earg}
Novamente, este é um argumento válido. A estrutura aqui é algo como:
	\begin{earg}
		\item[] $A$ ou $B$
		\item[] não-$A$
		\item[\texttherefore] $B$
	\end{earg}
Uma estrutura extraordinária! Eis outro exemplo:
	\begin{earg}
		\item[] Não é o caso de Jim ter estudado muito e atuado em muitas peças.
		\item[] Jim estudou muito.
		\item[\texttherefore] Jim não atuou em muitas peças.
	\end{earg}
Este argumento válido tem uma estrutura que poderíamos representar assim:
	\begin{earg}
		\item[] não-($A$ e $B$)
		\item[] $A$
		\item[\texttherefore] não-$B$
	\end{earg}
Esses exemplos ilustram uma ideia importante, que poderíamos descrever como \emph{validade em virtude da forma}. A validade dos argumentos que acabamos de considerar não tem relação significativa com os significados de expressões em português como Jenny está infeliz, Dipan é leitor ávido de Tolstói ou Jim atuou em muitas peças. Se tiver relação com significados, é com os significados de expressões como \emph{e}, \emph{ou}, \emph{não} e \emph{se ..., então ...}.
 
Em \crefrange{ch.TFL}{ch.NDTFL}, desenvolveremos uma
linguagem formal que nos permitirá simbolizar muitos argumentos de modo a
mostrar que eles são válidos em virtude de sua forma. Essa será a
linguagem da \emph{lógica proposicional verofuncional}, ou
LVF.\footnote{A lógica proposicional verofuncional também é frequentemente chamada de lógica (clássica) proposicional ou sentencial.}

\section{Validade por razões especiais}

Em \cref{s:FormalValidity}, introduzimos inicialmente a noção de
validade formal e a contrastamos com outros tipos de validade relacionados
aos tipos de contraexemplo que consideramos. Vale repetir que há muitos
argumentos que são válidos, mas não por razões relacionadas à sua forma.
Considere o seguinte exemplo:
	\begin{earg}
		\item[] Juanita é uma raposa fêmea.
		\item[\texttherefore] Juanita é uma raposa.
	\end{earg}
É impossível que a premissa seja verdadeira e a conclusão falsa, pois
raposa fêmea significa simplesmente raposa. Portanto, o argumento é
(conceitualmente) válido. A validade não é explicada pela forma do
argumento. Para perceber isso, podemos apresentar um argumento inválido
com a mesma forma, por exemplo:
	\begin{earg}
		\item[] Juanita é uma raposa fêmea.
		\item[\texttherefore] Juanita é uma catedral.
	\end{earg}
Considere também o argumento:
	\begin{earg}
		\item[] A escultura é verde por toda parte.
		\item[\texttherefore] A escultura não é vermelha por toda parte.
	\end{earg}
Novamente, parece não haver caso em que a premissa seja verdadeira e a
conclusão falsa, pois nada pode ser verde por toda parte e vermelho por
toda parte ao mesmo tempo. Portanto, o argumento é válido, mas eis um
argumento inválido com a mesma forma:
	\begin{earg}
		\item[] A escultura é verde por toda parte.
		\item[\texttherefore] A escultura não é brilhante por toda parte.
	\end{earg}
Este argumento é inválido, pois é possível ser verde por toda parte e
brilhante por toda parte. (Pode-se pintar a escultura com um elegante
verniz verde brilhante.) É plausível que a validade do primeiro argumento
esteja ligada à maneira como as cores (ou as palavras que designam cores)
interagem, mas, esteja isso correto ou não, não é simplesmente a
\emph{forma} do argumento, por si só, que o torna válido.

A lição importante pode ser enunciada assim. \emph{Na melhor das hipóteses, a LVF nos ajudará a compreender argumentos que são válidos em virtude de sua forma.}

\section{Sentenças atômicas}

Começamos a isolar a forma de um argumento, em
\cref{s:ValidityInVirtueOfForm}, substituindo \emph{subsentenças} de
sentenças por letras individuais. Assim, no primeiro exemplo desta seção,
está chovendo lá fora é uma subsentença de se está chovendo lá fora, então
Jenny está infeliz, e substituímos essa subsentença por $A$.

A linguagem artificial da LVF leva essa ideia às últimas consequências.
Começamos com algumas \emph{letras sentenciais}. Elas serão os blocos
básicos de construção a partir dos quais sentenças mais complexas serão
construídas (são as sentenças \emph{atômicas} da LVF). Usaremos letras
maiúsculas isoladas como letras sentenciais da LVF. Há apenas vinte e seis
letras no alfabeto, mas não há limite para o número de letras sentenciais
que podemos querer considerar. Ao acrescentar subscritos às letras,
obtemos novas letras sentenciais. Portanto, eis cinco letras sentenciais
diferentes da LVF:
	$$A, P, P_1, P_2, A_{234}$$
Usaremos letras sentenciais para representar, ou \emph{simbolizar},
certas sentenças em português. Para isso, fornecemos uma
\define{chave de simboliza\c{c}\~ao}, como a seguinte:
	\begin{ekey}
		\item[A] Está chovendo lá fora
		\item[C] Jenny está infeliz
	\end{ekey}
Ao fazer isso, não estamos fixando essa simbolização \emph{de uma vez
por todas}. Estamos apenas dizendo que, por enquanto, consideraremos a
letra sentencial da LVF, $A$, como simbolizando a sentença em português
Está chovendo lá fora, e a letra sentencial da LVF, $C$, como simbolizando
a sentença em português Jenny está infeliz. Mais tarde, quando estivermos
lidando com sentenças ou argumentos diferentes, poderemos fornecer uma
nova chave de simbolização, como esta:
	\begin{ekey}
		\item[A] Jenny é anarco-sindicalista
		\item[C] Dipan é leitor ávido de Tolstói
	\end{ekey}
É importante compreender que qualquer estrutura que uma sentença em
português possa ter se perde quando ela é simbolizada por uma letra
sentencial da LVF. Do ponto de vista da LVF, uma letra sentencial é apenas
uma letra. Ela pode ser usada para construir sentenças mais complexas,
mas não pode ser decomposta.

\newglossaryentry{sentence letter}
{
name=letra sentencial,
description={Uma letra usada para representar uma sentença básica na LVF}
}
\newglossaryentry{atomic sentence}
{
name=sentença atômica,
description={Uma expressão usada para representar uma sentença básica: uma letra sentencial na LVF ou um símbolo de predicado seguido de nomes na LPO}
}
\newglossaryentry{symbolization key}
{
name=chave de simbolização,
description={Uma lista que mostra quais sentenças em português são representadas por quais \glspl{sentence letter} na LVF}
}
\chapter{Conectivos}
\label{s:TFLConnectives}

No capítulo anterior, consideramos a simbolização de sentenças em
português relativamente básicas usando letras sentenciais da LVF. Isso
nos deixa querendo lidar com as expressões em português \emph{e},
\emph{ou}, \emph{não} e assim por diante. Elas são \emph{conectivos}:
podem ser usadas para formar novas sentenças a partir de sentenças
antigas. Na LVF, usaremos conectivos lógicos para construir sentenças
complexas a partir de componentes atômicos. Há cinco conectivos lógicos
na LVF. A tabela a seguir os resume, e eles são explicados ao longo
desta seção.

\newglossaryentry{connective}
{
name=conectivo,
description={Um operador lógico na LVF usado para combinar \glspl{sentence letter} em sentenças maiores}
}
	\begin{table}[h]
	\centering
	\begin{tabular}{l l l}
	
	\textbf{símbolo}&\textbf{como é chamado}&\textbf{significado aproximado}\\
	\hline
	\enot&negação&Não é o caso que$\ldots$\\
	\eand&conjunção&Ambos$\ldots$\ e $\ldots$\\
	\eor&disjunção&Ou$\ldots$\ ou $\ldots$\\
	\eif&condicional&Se $\ldots$\ então $\ldots$\\
	\eiff&bicondicional&$\ldots$\ se e somente se $\ldots$\\
	
	\end{tabular}
	\end{table}

Esses não são os únicos conectivos da língua portuguesa que nos
interessam. Outros são, por exemplo, \emph{a menos que}, \emph{nem
\dots{} nem \dots} e \emph{porque}. Veremos que os dois primeiros podem
ser expressos pelos conectivos que discutiremos, enquanto o último não
pode. \emph{Porque}, ao contrário dos demais, não é
\emph{verofuncional}.

\section{Negação}

Considere como poderíamos simbolizar estas sentenças:
	\begin{enumerate}
	\item\label{not1} Mary está em Barcelona.
	\item\label{not2} Não é o caso de Mary estar em Barcelona.
	\item\label{not3} Mary não está em Barcelona.
	\end{enumerate}
Para simbolizar \cref*{not1}, precisaremos de uma letra sentencial.
Poderíamos oferecer esta chave de simbolização:
	\begin{ekey}
		\item[B] Mary está em Barcelona.
	\end{ekey}
Como \cref*{not2} está obviamente relacionada a \cref*{not1}, não
quereremos simbolizá-la com uma letra sentencial completamente
diferente. Grosso modo, \cref*{not2} significa algo como
\emph{não é o caso de $B$}. Para simbolizar isso, precisamos de um
símbolo para a negação. Usaremos \enot. Agora podemos simbolizar
\cref*{not2} como $\enot B$.

\Cref*{not3} também contém a palavra \emph{não}, e é obviamente
equivalente a \cref*{not2}. Assim, também podemos simbolizá-la como
$\enot B$.
\factoidbox{
Uma sentença pode ser simbolizada como $\enot\metav{A}$ se puder ser
parafraseada em português como \emph{não é o caso que\ldots}.
}
Será útil oferecer mais alguns exemplos:
	\begin{enumerate}
		\item\label{not4} O dispositivo pode ser substituído.
		\item\label{not5} O dispositivo é insubstituível.
		\item\label{not5b} O dispositivo não é insubstituível.
	\end{enumerate}
Usaremos a seguinte chave de representação:
	\begin{ekey}
		\item[R] O dispositivo pode ser substituído.
	\end{ekey}
Agora \cref*{not4} pode ser simbolizada por $R$. Prosseguindo para
\cref*{not5}: dizer que o dispositivo é insubstituível significa que não
é o caso de o dispositivo poder ser substituído. Portanto, embora
\cref*{not5} não contenha a palavra \emph{não}, nós a simbolizaremos
assim: $\enot R$.

\Cref*{not5b} pode ser parafraseada como \emph{não é o caso de o
dispositivo ser insubstituível}. Isso pode novamente ser parafraseado
como \emph{não é o caso de não ser o caso de o dispositivo poder ser
substituído}. Assim, poderíamos simbolizar essa sentença em português
com a sentença da LVF $\enot\enot R$.

Mas é preciso algum cuidado ao lidar com negações. Considere:
	\begin{enumerate}
		\item\label{not6} Jane está feliz.
		\item\label{not7} Jane está infeliz.
	\end{enumerate}
Se deixarmos a sentença da LVF $H$ simbolizar \emph{Jane está feliz},
poderemos simbolizar \cref*{not6} como $H$. Entretanto, seria um erro
simbolizar \cref*{not7} por $\enot{H}$. Se Jane está infeliz, então ela
não está feliz; mas \cref*{not7} não significa o mesmo que \emph{não é o
caso de Jane estar feliz}. Jane pode não estar nem feliz nem infeliz;
pode estar em um estado de indiferença absoluta. Portanto, para
simbolizar \cref*{not7}, precisaríamos de uma nova letra sentencial da
LVF.

\newglossaryentry{negation}
{
name=negação,
description={O símbolo \enot, usado para representar palavras e frases que funcionam como a palavra portuguesa não}
}

\section{Conjunção}
\label{s:ConnectiveConjunction}

Considere estas sentenças:
	\begin{enumerate}
		\item\label{and1}Adam é atlético.
		\item\label{and2}Barbara é atlética.
		\item\label{and3}Adam é atlético, e Barbara também é atlética.
	\end{enumerate}
Precisaremos de letras sentenciais separadas da LVF para simbolizar
\cref*{and1,and2}; por exemplo:
	\begin{ekey}
		\item[A] Adam é atlético.
		\item[B] Barbara é atlética.
	\end{ekey}
\Cref*{and1} pode agora ser simbolizada como $A$, e \cref*{and2} pode
ser simbolizada como $B$. \Cref*{and3} diz aproximadamente $A$ e $B$.
Precisamos de outro símbolo para lidar com \emph{e}. Usaremos
\emph{\eand}. Assim, vamos simbolizá-la como $(A\eand B)$. Esse
conectivo é chamado de \define{conjun\c{c}\~ao}. Também dizemos que
$A$ e $B$ são os dois \define{conjuntos} da conjunção $(A \eand B)$.

\newglossaryentry{conjunction}
{
name=conjunção,
description={O símbolo \eand, usado para representar palavras e frases que funcionam como a palavra portuguesa e; ou uma sentença formada usando esse símbolo}
}

\newglossaryentry{conjunct}
{
name=conjunto,
description={Uma sentença ligada a outra por uma \gls{conjunction}}
}


Observe que não fazemos nenhuma tentativa de simbolizar a palavra
\emph{também} em \cref*{and3}. Palavras como \emph{ambos} e
\emph{também} servem para chamar nossa atenção para o fato de que duas
coisas estão sendo ligadas por uma conjunção. Talvez elas afetem a ênfase de uma
sentença, mas não vamos (e não podemos) simbolizar tais coisas na LVF.

Mais alguns exemplos tornarão esse ponto claro:
	\begin{enumerate}
		\item\label{and4}Barbara é atlética e enérgica.
		\item\label{and5}Barbara e Adam são ambos atléticos.
		\item\label{and6}Embora Barbara seja enérgica, ela não é atlética.
	\item\label{and7}Adam é atlético, mas Barbara é mais atlética do que ele.
	\end{enumerate}
\Cref*{and4} é obviamente uma conjunção. A sentença diz duas coisas
sobre Barbara. Em português, é permitido referir-se a Barbara apenas
uma vez. Pode ser tentador pensar que precisamos simbolizar
\cref*{and4} com algo como $B$ e enérgica. Isso seria um erro. Depois
que simbolizamos parte de uma sentença como $B$, toda estrutura
adicional se perde, pois $B$ é uma letra sentencial da LVF. Inversamente,
\emph{enérgica} nem sequer é uma sentença em português. O que buscamos
é algo como $B$ e Barbara é enérgica. Portanto, precisamos acrescentar
outra letra sentencial à chave de simbolização. Seja $E$ a letra que
simboliza Barbara é enérgica. Agora a sentença inteira pode ser
simbolizada como $(B\eand E)$.

\Cref*{and5} diz uma coisa sobre dois sujeitos diferentes. Diz tanto de
Barbara quanto de Adam que são atléticos, embora em português usemos a
palavra \emph{atlético} apenas uma vez. A sentença pode ser parafraseada
como Barbara é atlética, e Adam é atlético. Podemos simbolizá-la na LVF
como $(B\eand A)$, usando a mesma chave de simbolização que estamos
usando.

\Cref{and6} é um pouco mais complicada. A palavra \emph{embora}
estabelece um contraste entre a primeira parte da sentença e a segunda.
Ainda assim, a sentença nos diz tanto que Barbara é enérgica quanto que
ela não é atlética. Para tornar cada um dos conjuntos uma letra
sentencial, precisamos substituir \emph{ela} por \emph{Barbara}. Assim,
podemos reescrever \cref*{and6} como: \emph{Ambos} Barbara é enérgica
\emph{e} Barbara não é atlética. O segundo conjunto contém uma negação,
portanto parafraseamos ainda mais: \emph{Ambos} Barbara é enérgica
\emph{e não é o caso de} Barbara ser atlética. Agora podemos simbolizar
isso com a sentença da LVF $(E\eand\enot B)$. Observe que perdemos todo
tipo de nuance nessa simbolização. Há uma diferença distinta de tom
entre \cref*{and6} e \emph{Ambos Barbara é enérgica e não é o caso de
Barbara ser atlética}. A LVF não preserva (e não pode preservar) essas
nuances.

\Cref{and7} apresenta questões semelhantes. A palavra \emph{mas} sugere
um contraste ou uma diferença, mas isso não é algo com que a LVF possa
lidar. Tudo que podemos fazer é parafrasear a sentença como
\emph{Ambos} Adam é atlético, \emph{e} Barbara é mais atlética do que
Adam. (Observe que, mais uma vez, substituímos o pronome \emph{ele} por
\emph{Adam}.) Como devemos lidar com o segundo conjunto? Já temos a
letra sentencial $A$, usada para simbolizar Adam é atlético, e a sentença
$B$, usada para simbolizar Barbara é atlética; mas nenhuma delas diz
respeito à atléticidade relativa das duas pessoas. Portanto, para
simbolizar a sentença inteira, precisamos de uma nova letra sentencial.
Seja a sentença da LVF $R$ a letra que simboliza a sentença em português
Barbara é mais atlética do que Adam. Agora podemos simbolizar
\cref*{and7} por $(A \eand R)$.
\factoidbox{Uma sentença pode ser simbolizada como
$(\metav{A}\eand\metav{B})$ se puder ser parafraseada em português como
\emph{Ambos\ldots, e\ldots}, ou como \emph{\ldots, mas\ldots}, ou como
\emph{embora \ldots, \ldots}.}

Observamos acima que o contraste sugerido por \emph{mas} não pode ser
capturado na LVF e que simplesmente o ignoramos. Um fenômeno que não
pode ser simplesmente ignorado é a ordem temporal. Considere, por
exemplo:
\begin{enumerate}
	\item\label{stoodfirst} Harry se levantou e objetou à proposta.
	\item\label{stoodsecond} Harry objetou à proposta e se levantou.
\end{enumerate}
Se Harry se levantou \emph{depois} de objetar, \cref*{stoodsecond} é
verdadeira, mas \cref*{stoodfirst} é falsa: o uso de \emph{e} aqui é
\emph{assimétrico}. O símbolo \eand da LVF, contudo, é sempre simétrico
(ou \emph{comutativo}, como dizem os lógicos). A LVF não pode lidar com um
\emph{e} assimétrico. Vamos supor, em todos os nossos exemplos e
exercícios, que \emph{e} é simétrico.\footnote{Sobre a conjunção
simétrica e assimétrica na linguística, veja, por exemplo, Robin Lakoff,
\textit{If's, and's, and but's about conjunction}, em: C. J. Fillmore e
D. T. Langendoen (orgs.), \textit{Studies in Linguistic Semantics},
Holt, Rinehart \& Winston, 1971, e Susan Schmerling, \textit{Asymmetric
conjunction and rules of conversation}, em: P. Cole e J. L. Morgan
(orgs.), \textit{Speech Acts}, Brill, 1975, p.~211--31.}

Você pode estar se perguntando por que colocamos parênteses em torno das
conjunções. O motivo pode ser esclarecido pensando em como a negação
interage com a conjunção. Considere:
	\begin{enumerate}
		\item\label{negcon1}Não é o caso de você receber sopa e salada.
		\item\label{negcon2}Você não receberá sopa, mas receberá salada.
	\end{enumerate}
\Cref*{negcon1} pode ser parafraseada como \emph{não é o caso de: você
receberá sopa e você receberá salada}. Usando esta chave de simbolização:
	\begin{ekey}
		\item[\ensuremath{S_1}] Você receberá sopa.
		\item[\ensuremath{S_2}] Você receberá salada.
	\end{ekey}
simbolizaríamos \emph{você receberá sopa e você receberá salada} como
$(S_1 \eand S_2)$. Para simbolizar \cref*{negcon1}, então, simplesmente
negamos a sentença inteira, assim: $\enot (S_1 \eand S_2)$.

\Cref*{negcon2} é uma conjunção: você \emph{não receberá} sopa e
\emph{receberá} salada. \emph{Você não receberá sopa} é simbolizada por
$\enot S_1$. Portanto, para simbolizar a própria \cref*{negcon2},
oferecemos $(\enot S_1 \eand S_2)$.

Essas sentenças em português são muito diferentes, e suas simbolizações
diferem de acordo. Em uma delas, a conjunção inteira é negada. Na outra,
apenas um conjunto é negado. Os parênteses nos ajudam a acompanhar
coisas como o \emph{escopo} da negação.

\section{Disjunção}

Considere estas sentenças:
	\begin{enumerate}
		\item\label{or1}Ou Fatima jogará videogame, ou assistirá a filmes.
		\item\label{or2}Ou Fatima ou Omar jogará videogame.
	\end{enumerate}
Para essas sentenças, podemos usar esta chave de simbolização:
	\begin{ekey}
		\item[F] Fatima jogará videogame
		\item[O] Omar jogará videogame
		\item[M] Fatima assistirá a filmes
	\end{ekey}
Entretanto, novamente precisaremos introduzir um novo símbolo.
\Cref*{or1} é simbolizada por $(F \eor M)$. O conectivo é chamado de
\define{disjun\c{c}\~ao}. Também dizemos que $F$ e $M$ são os
\define{disjuntos} da disjunção $(F \eor M)$.

\newglossaryentry{disjunction}
{
name=disjunção,
description={O conectivo \eor, usado para representar palavras e frases que funcionam como a palavra portuguesa ou em seu sentido inclusivo; ou uma sentença formada usando esse conectivo}
}

\newglossaryentry{disjunct}
{
name=disjunto,
description={Uma sentença ligada a outra por uma \gls{disjunction}}
}

\Cref*{or2} é apenas um pouco mais complicada. Há dois sujeitos, mas a
sentença em português apresenta o verbo apenas uma vez. Entretanto,
podemos parafrasear \cref*{or2} como Ou Fatima jogará videogame, ou Omar
jogará videogame. Agora podemos obviamente simbolizá-la novamente por
$(F \eor O)$.
	\factoidbox{
		Uma sentença pode ser simbolizada como $(\metav{A}\eor\metav{B})$ se puder ser parafraseada em português como \emph{Ou\ldots, ou\ldots}.
	}
Às vezes, em português, a palavra \emph{ou} é usada de uma maneira que
exclui a possibilidade de ambos os disjuntos serem verdadeiros. Isso é
chamado de \define{ou exclusivo}. Um \emph{ou exclusivo} é claramente
pretendido quando um cardápio diz: \emph{As entradas vêm com sopa ou
salada}: você pode escolher sopa; pode escolher salada; mas, se quiser
\emph{ambas} a sopa \emph{e} a salada, terá de pagar a mais.

Em outras ocasiões, a palavra \emph{ou} permite a possibilidade de ambos
os disjuntos serem verdadeiros. Esse provavelmente é o caso de
\cref{or2}, acima. Fatima pode jogar videogame sozinha, Omar pode jogar
videogame sozinho, ou ambos podem jogar. \Cref*{or2} apenas diz que
\emph{pelo menos} um deles joga videogame. Isso é um \define{ou inclusivo}. O símbolo \eor da LVF sempre simboliza um \emph{ou inclusivo}.

Também será útil ver como a negação interage com a disjunção. Considere:
	\begin{enumerate}
		\item\label{or3}Ou você não receberá sopa, ou não receberá salada.
		\item\label{or4}Você não receberá nem sopa nem salada.
		\item\label{or.xor}Você receberá sopa ou salada, mas não ambas.
	\end{enumerate}
Usando a mesma chave de simbolização de antes, \cref*{or3} pode ser
parafraseada desta maneira: \emph{Ou não é o caso de você receber sopa,
ou não é o caso de você receber salada}. Para simbolizar isso na LVF,
precisamos tanto da disjunção quanto da negação. \emph{Não é o caso de
você receber sopa} é simbolizada por $\enot S_1$. \emph{Não é o caso de
você receber salada} é simbolizada por $\enot S_2$. Portanto,
\cref*{or3} é simbolizada por $(\enot S_1 \eor \enot S_2)$.

\Cref{or4} também requer negação. Ela pode ser parafraseada como
\emph{não é o caso de: você receber sopa ou você receber salada}. Como
isso nega a disjunção inteira, simbolizamos \cref*{or4} por
$\enot (S_1 \eor S_2)$.

\Cref{or.xor} é um \emph{ou exclusivo}. Podemos dividir a sentença em
duas partes. A primeira parte diz que você recebe uma coisa ou a outra.
Nós a simbolizamos como $(S_1 \eor S_2)$. A segunda parte diz que você
não recebe ambas. Podemos parafrasear isso como: \emph{não é o caso de
você receber sopa e receber salada}. Usando tanto a negação quanto a
conjunção, simbolizamos isso por $\enot(S_1 \eand S_2)$. Agora só
precisamos juntar as duas partes. Como vimos acima, \emph{mas} pode
costumeiramente ser simbolizado por $\eand$. Assim, \cref*{or.xor} pode
ser simbolizada como $((S_1 \eor S_2) \eand \enot(S_1 \eand S_2))$.

Este último exemplo mostra algo importante. Embora o símbolo \eor da LVF
sempre simbolize \emph{ou inclusivo}, podemos simbolizar um \emph{ou
exclusivo} na LVF. Basta usar também alguns outros símbolos.

\section{Condicional}
Considere estas sentenças:
	\begin{enumerate}
		\item\label{if1}Se Jean está em Paris, então Jean está na França.
		\item\label{if2}Jean está na França somente se Jean está em Paris.
	\end{enumerate}
Usaremos a seguinte chave de simbolização:
	\begin{ekey}
		\item[P] Jean está em Paris
		\item[F] Jean está na França
	\end{ekey}
\Cref*{if1} tem aproximadamente esta forma: \emph{se $P$, então
$F$}. Usaremos o símbolo \eif para simbolizar essa estrutura de
\emph{se\ldots, então\ldots}. Portanto, simbolizamos \cref*{if1} por
$(P\eif F)$. O conectivo é chamado de \define{condicional}. Aqui, $P$ é
chamado de \define{antecedente} do condicional $(P \eif F)$, e $F$ é
chamado de \define{consequente}.

\newglossaryentry{conditional}
{
name=condicional,
description={O símbolo \eif, usado para representar palavras e frases que funcionam como a expressão portuguesa se \dots{} então \dots{}; uma sentença formada usando esse símbolo}
}

\newglossaryentry{antecedent}
{
name=antecedente,
description={A sentença do lado esquerdo de um \gls{conditional}}
}


\newglossaryentry{consequent}
{
name=consequente,
description={A sentença do lado direito de um \gls{conditional}}
}

\Cref{if2} também é um condicional. Como a palavra \emph{se} aparece na
segunda metade da sentença, pode ser tentador simbolizá-la da mesma
maneira que \cref*{if1}. Isso seria um erro. Nosso conhecimento de
geografia nos diz que \cref*{if1} é inequivocamente verdadeira: não há
como Jean estar em Paris sem estar na França. Mas \cref*{if2} não é tão
simples: se Jean estivesse em Dieppe, Lyon ou Toulouse, estaria na
França sem estar em Paris, tornando \cref*{if2} falsa. Como a geografia,
por si só, determina a verdade de \cref*{if1}, enquanto planos de viagem,
por exemplo, são necessários para saber a verdade de \cref*{if2}, elas
devem significar coisas diferentes.

Na verdade, \cref*{if2} pode ser parafraseada como \emph{Se Jean está na
França, então Jean está em Paris}. Portanto, podemos simbolizá-la por
$(F \eif P)$.
	\factoidbox{
		Uma sentença pode ser simbolizada como $(\metav{A} \eif \metav{B})$ se puder ser parafraseada em português como \emph{Se A, então B} ou \emph{A somente se B}.
	}
\noindent Na verdade, o condicional pode representar muitas expressões
em português. Considere:
	\begin{enumerate}
		\item\label{ifnec1}Para Jean estar em Paris, é necessário que Jean esteja na França.
		\item\label{ifnec2}É uma condição necessária para Jean estar em Paris que ela esteja na França.
		\item\label{ifsuf1}Para Jean estar na França, é suficiente que Jean esteja em Paris.
		\item\label{ifsuf2}É uma condição suficiente para Jean estar na França que ela esteja em Paris.
	\end{enumerate}
Pensando bem, todas essas quatro sentenças significam o mesmo que
\emph{Se Jean está em Paris, então Jean está na França}. Portanto, todas
podem ser simbolizadas por $(P \eif F)$.

É importante ter em mente que o conectivo \eif nos diz apenas que, se o
antecedente é verdadeiro, então o consequente é verdadeiro. Ele nada diz
sobre uma conexão \emph{causal} entre dois eventos, por exemplo. Na
verdade, perdemos uma quantidade enorme de informação quando usamos
\eif para simbolizar condicionais em português. Retornaremos a isso em
\cref{s:IndicativeSubjunctive} e \cref{s:ParadoxesOfMaterialConditional}.

\section{Bicondicional}
Considere estas sentenças:
	\begin{enumerate}
		\item\label{iff1}Laika é uma cadela somente se ela é mamífera.
		\item\label{iff2}Laika é uma cadela se ela é mamífera.
		\item\label{iff3}Laika é uma cadela se e somente se ela é mamífera.
	\end{enumerate}
Usaremos a seguinte chave de simbolização:
	\begin{ekey}
		\item[D] Laika é uma cadela
		\item[M] Laika é mamífera
	\end{ekey}
Pelas razões discutidas acima, \cref*{iff1} pode ser simbolizada por
$(D \eif M)$.


\Cref*{iff3} diz algo mais forte que \cref*{iff1} ou \cref*{iff2}.
Pode ser parafraseada como \emph{Laika é uma cadela se Laika é mamífera,
e Laika é uma cadela somente se Laika é mamífera}. Isso é simplesmente
a conjunção de \cref*{iff1,iff2}. Portanto, podemos simbolizá-la como
$(D \eif M) \eand (M \eif D)$. Chamamos isso de
\define{bicondicional}, porque equivale a afirmar as duas direções do
condicional.

\newglossaryentry{biconditional}
{
name=bicondicional,
description={O símbolo \eiff, usado para representar palavras e frases que funcionam como a expressão portuguesa se e somente se; ou uma sentença formada usando esse conectivo}
}

Poderíamos tratar todo bicondicional dessa maneira. Assim, do mesmo modo
que não precisamos de um novo símbolo da LVF para lidar com o
\emph{ou exclusivo}, também não precisamos realmente de um novo símbolo
da LVF para lidar com bicondicionais. Como o bicondicional ocorre com
muita frequência, contudo, usaremos o símbolo \eiff para ele. Então
podemos simbolizar \cref*{iff3} com a sentença da LVF $(D \eiff M)$.

A expressão \emph{se e somente se} ocorre com muita frequência,
especialmente em filosofia, matemática e lógica. Por brevidade, podemos
abreviá-la pela palavra mais concisa \emph{iff}. Seguiremos essa prática.
Assim, \emph{if}, com apenas um \emph{f}, é o condicional em inglês, mas
\emph{iff}, com dois \emph{f}s, é o bicondicional em inglês. Munidos
disso, podemos dizer:
	\factoidbox{
		Uma sentença pode ser simbolizada como $(\metav{A} \eiff \metav{B})$ se puder ser parafraseada em português como \emph{A \ifeff{} B}; isto é, como \emph{A se e somente se B}.
	}
Uma palavra de cautela. Falantes comuns de português frequentemente
usam \emph{se\ldots, então\ldots} quando na verdade querem dizer algo
mais próximo de \emph{se e somente se}. Talvez seus pais tenham dito a
você, quando era criança: \emph{se você não comer seus vegetais, não
ganhará sobremesa}. Suponha que você tenha comido seus vegetais, mas que
seus pais tenham se recusado a lhe dar sobremesa, alegando que estavam
comprometidos apenas com o \emph{condicional} (aproximadamente, \emph{se
você receber sobremesa, então terá comido seus vegetais}), e não com o
bicondicional (aproximadamente, \emph{você recebe sobremesa \ifeff{}
come seus vegetais}). Bem, uma birra seria justificadamente inevitável.
Portanto, fique atento a isso ao interpretar as pessoas; mas, em sua
própria escrita, certifique-se de usar o bicondicional \ifeff{} quando
quiser dizer isso.

\section{A menos que}
Agora introduzimos todos os conectivos da LVF. Podemos usá-los juntos
para simbolizar muitos tipos de sentenças. Um caso especialmente difícil
ocorre quando usamos o conectivo da língua portuguesa \emph{a menos que}:

\begin{enumerate}
\item\label{unless1}A menos que você vista uma jaqueta, ficará resfriado.
\item\label{unless2}Você ficará resfriado a menos que vista uma jaqueta.
\end{enumerate}
Essas duas sentenças são claramente equivalentes. Para simbolizá-las,
usaremos a chave de simbolização:
	\begin{ekey}
		\item[J] Você vestirá uma jaqueta.
		\item[D] Você ficará resfriado.
	\end{ekey}
Ambas as sentenças significam que, se você não vestir uma jaqueta, ficará
resfriado. Com isso em mente, poderíamos simbolizá-las como
$(\enot J \eif D)$.

Da mesma forma, ambas as sentenças significam que, se você não ficar
resfriado, então necessariamente vestiu uma jaqueta. Com isso em mente,
poderíamos simbolizá-las como $(\enot D \eif J)$.

Da mesma forma, ambas as sentenças significam que ou você vestirá uma
jaqueta ou ficará resfriado. Com isso em mente, poderíamos simbolizá-las
como $(J \eor D)$.

As três são simbolizações corretas. De fato, em
\cref{s:SemanticConcepts}, veremos que as três simbolizações são
equivalentes na LVF.
% TODO: pode ser útil fazer aqui uma referência explícita ao exercício
% 11.F.3, pois o ponto não é discutido no texto principal
	\factoidbox{
		Se uma sentença pode ser parafraseada como \emph{A menos que $A$, $B$}, então pode ser simbolizada como $(\enot\metav{A}\eif\metav{B})$
		ou $(\metav{A}\eor\metav{B})$.
	}
Novamente, porém, há uma pequena complicação. \emph{A menos que} pode
ser simbolizado como um condicional; mas, como dissemos acima, as pessoas
frequentemente usam o condicional isoladamente quando querem dizer algo
mais parecido com o bicondicional. Da mesma forma, \emph{a menos que}
pode ser simbolizado como uma disjunção; mas há dois tipos de disjunção
(exclusiva e inclusiva). Portanto, não será surpresa descobrir que
falantes comuns frequentemente usam \emph{a menos que} para significar
algo mais parecido com o bicondicional ou com a disjunção exclusiva.
Suponha que alguém diga: \emph{Vou correr a menos que chova}. Essa pessoa
provavelmente quer dizer algo como \emph{Vou correr \ifeff{} não chover}
(isto é, o bicondicional) ou \emph{ou vou correr ou choverá, mas não
ambos} (isto é, a disjunção exclusiva). Novamente: esteja atento a isso
ao interpretar o que outras pessoas disseram, mas seja preciso em sua
escrita.

\practiceproblems
\solutions
\problempart Usando a chave de simbolização fornecida, simbolize cada sentença em português na LVF.\label{pr.monkeysuits}
	\begin{ekey}
		\item[M] Essas criaturas são homens de terno
		\item[C] Essas criaturas são chimpanzés
		\item[G] Essas criaturas são gorilas
	\end{ekey}
\begin{compactlist}
\item Essas criaturas não são homens de terno.
\item Essas criaturas são homens de terno, ou não são.
\item Essas criaturas são ou gorilas ou chimpanzés.
\item Essas criaturas não são nem gorilas nem chimpanzés.
\item Se essas criaturas são chimpanzés, então não são nem gorilas nem homens de terno.
\item A menos que essas criaturas sejam homens de terno, elas são ou chimpanzés ou gorilas.
\end{compactlist}

\problempart Usando a chave de simbolização fornecida, simbolize cada sentença em português na LVF.
\begin{ekey}
\item[A] O Senhor Ace foi assassinado
\item[B] O mordomo foi o responsável
\item[C] O cozinheiro foi o responsável
\item[D] A Duquesa está mentindo
\item[E] O Senhor Edge foi assassinado
\item[F] A arma do crime foi uma frigideira
\end{ekey}
\begin{compactlist}
\item Ou o Senhor Ace ou o Senhor Edge foi assassinado.
\item Se o Senhor Ace foi assassinado, então o cozinheiro foi o responsável.
\item Se o Senhor Edge foi assassinado, então o cozinheiro não foi o responsável.
\item Ou o mordomo foi o responsável, ou a Duquesa está mentindo.
\item O cozinheiro foi o responsável somente se a Duquesa está mentindo.
\item Se a arma do crime foi uma frigideira, então o culpado deve ter sido o cozinheiro.
\item Se a arma do crime não foi uma frigideira, então o culpado foi o cozinheiro ou o mordomo.
\item O Senhor Ace foi assassinado se e somente se o Senhor Edge não foi assassinado.
\item A Duquesa está mentindo, a menos que tenha sido o Senhor Edge quem foi assassinado.
\item Se o Senhor Ace foi assassinado, ele foi morto com uma frigideira.
\item Como o cozinheiro foi o responsável, o mordomo não foi.
\item É claro que a Duquesa está mentindo!
\end{compactlist}

\solutions

\problempart Usando a chave de simbolização fornecida, simbolize cada sentença em português na LVF.\label{pr.avacareer}
	\begin{ekey}
		\item[\ensuremath{E_1}] Ava é eletricista
		\item[\ensuremath{E_2}] Harrison é eletricista
		\item[\ensuremath{F_1}] Ava é bombeira
		\item[\ensuremath{F_2}] Harrison é bombeiro
		\item[\ensuremath{S_1}] Ava está satisfeita com sua carreira
		\item[\ensuremath{S_2}] Harrison está satisfeito com sua carreira
	\end{ekey}
\begin{compactlist}
\item Ava e Harrison são ambos eletricistas.
\item Se Ava é bombeira, então está satisfeita com sua carreira.
\item Ava é bombeira, a menos que seja eletricista.
\item Harrison é um eletricista insatisfeito.
\item Nem Ava nem Harrison é eletricista.
\item Ava e Harrison são ambos eletricistas, mas nenhum dos dois considera isso satisfatório.
\item Harrison está satisfeito somente se é bombeiro.
\item Se Ava não é eletricista, Harrison também não é, mas, se ela é, ele também é.
\item Ava está satisfeita com sua carreira se e somente se Harrison não está satisfeito com a sua.
\item Se Harrison é tanto eletricista quanto bombeiro, então deve estar satisfeito com seu trabalho.
\item Não pode ser que Harrison seja tanto eletricista quanto bombeiro.
\item Harrison e Ava são ambos bombeiros se e somente se nenhum dos dois é eletricista.
\end{compactlist}

\problempart
Usando a chave de simbolização fornecida, traduza cada sentença em português para a LVF.
\label{pr.jazzinstruments}
\begin{ekey}
\item[\ensuremath{J_1}] John Coltrane tocou saxofone tenor
\item[\ensuremath{J_2}] John Coltrane tocou saxofone soprano
\item[\ensuremath{J_3}] John Coltrane tocou tuba
\item[\ensuremath{M_1}] Miles Davis tocou trompete
\item[\ensuremath{M_2}] Miles Davis tocou tuba
\end{ekey}

\begin{compactlist}
\item John Coltrane tocou saxofone tenor e soprano. %{\color{red} $J_1 \eand J_2$} \vspace{1ex}
\item Nem Miles Davis nem John Coltrane tocaram tuba. %{\color{red} $\enot(M_2 \eor J_3)$ or $\enot M_2 \eand \enot J_3$} \vspace{1ex}
\item John Coltrane não tocou saxofone tenor e tuba ao mesmo tempo. %{\color{red} $\enot(J_1 \eand J_3)$ or $\enot J_1 \eor \enotJ_3$} \vspace{1ex}
\item John Coltrane não tocou saxofone tenor a menos que também tenha tocado saxofone soprano. %{\color{red} $\enot J_1 \eor J_2$} \vspace{1ex}
\item John Coltrane não tocou tuba, mas Miles Davis tocou. %{\color{red} $\enotJ_3 \eand M_2$} \vspace{1ex}
\item Miles Davis tocou trompete somente se também tocou tuba. %{\color{red} $M_1 \eiff M_2$} \vspace{1ex}
\item Se Miles Davis tocou trompete, então John Coltrane tocou pelo menos um destes três instrumentos: saxofone tenor, saxofone soprano ou tuba. %{\color{red} $M_1 \eif (J_1 \eor (J_2 \eor J_3))&} \vspace{1ex}
\item Se John Coltrane tocou tuba, então Miles Davis não tocou nem trompete nem tuba. %{\color{red} $J_3 \eif \enot(M_1 \eor M_2)$ or $J_3 \eif (\enot M_1 \eand \enot M_2)$  } \vspace{1ex}
\item Miles Davis e John Coltrane tocaram tuba ambos se e somente se Coltrane não tocou saxofone tenor e Miles Davis não tocou trompete. %{\color{red} $(J_3 \eand M_2) \eiff \enotJ_1 & \enot M_1)$ or $(J_3 \eand M_2) \eiff \enot (J_1 \eor M_1)$} \vspace{1ex}
\end{compactlist}

\solutions
\problempart
\label{pr.spies}
Dê uma chave de simbolização e simbolize as seguintes sentenças em português na LVF.
\begin{compactlist}
\item Alice e Bob são ambos espiões.
\item Se Alice ou Bob é espião, então o código foi quebrado.
\item Se nem Alice nem Bob é espião, então o código permanece intacto.
\item A embaixada alemã ficará em alvoroço, a menos que alguém tenha quebrado o código.
\item Ou o código foi quebrado ou não foi, mas a embaixada alemã ficará em alvoroço de todo modo.
\item Alice ou Bob é espião, mas não ambos.
\end{compactlist}

\solutions
\problempart Dê uma chave de simbolização e simbolize as seguintes sentenças em português na LVF.
\begin{compactlist}
\item Se houver comida nas Terras do Reino, então Rafiki falará sobre bananas amassadas.
\item Rafiki falará sobre bananas amassadas a menos que Simba esteja vivo.
\item Rafiki falará sobre bananas amassadas ou não falará, mas há comida nas Terras do Reino de todo modo.
\item Scar permanecerá rei se e somente se houver comida nas Terras do Reino.
\item Se Simba estiver vivo, então Scar não permanecerá rei.
\end{compactlist}

\problempart
Para cada argumento, escreva uma chave de simbolização e simbolize todas as sentenças do argumento na LVF.
\begin{compactlist}
\item Se Dorothy toca piano pela manhã, então Roger acorda mal-humorado. Dorothy toca piano pela manhã a menos que esteja distraída. Portanto, se Roger não acordar mal-humorado, então Dorothy deve estar distraída.
\item Na terça-feira, choverá ou nevará. Se chover, Neville ficará triste. Se nevar, Neville ficará com frio. Portanto, na terça-feira, Neville ficará triste ou com frio.
\item Se Zoog se lembrou de fazer suas tarefas, então as coisas estão limpas, mas não arrumadas. Se ele esqueceu, então as coisas estão arrumadas, mas não limpas. Portanto, as coisas estão arrumadas ou limpas, mas não ambas.
\end{compactlist}

\problempart
Para cada argumento, escreva uma chave de simbolização e simbolize o argumento tanto quanto possível na LVF. A parte do trecho em itálico fornece contexto para o argumento e não precisa ser simbolizada.
\begin{compactlist}
\item Vai chover em breve. Sei disso porque minha perna está doendo, e minha perna dói se vai chover.

\item  \emph{O Homem-Aranha tenta descobrir o plano do vilão.} Se o Doutor Octopus obtiver o urânio, chantageará a cidade. Tenho certeza disso porque, se o Doutor Octopus obtiver o urânio, poderá fabricar uma bomba suja e, se puder fabricar uma bomba suja, chantageará a cidade.

\item \emph{Um ocidental tenta prever as políticas do governo chinês.} Se o governo chinês não conseguir resolver a escassez de água em Pequim, terá de transferir a capital. Eles não querem transferir a capital. Portanto, precisam resolver a escassez de água. Mas a única maneira de resolver a escassez de água é desviar quase toda a água do rio Yangtzé para o norte. Portanto, o governo chinês seguirá o projeto de desviar água do sul para o norte.

\end{compactlist}


\problempart
Simbolizamos um \emph{ou exclusivo} usando \eor, \eand e \enot. Como você poderia simbolizar um \emph{ou exclusivo} usando apenas dois conectivos? Há alguma maneira de simbolizar um \emph{ou exclusivo} usando apenas um conectivo?

\chapter{Sentences of TFL}\label{s:TFLSentences}
The sentence `either apples are red, or berries are blue' is a sentence of English, and the sentence `$(A\eor B)$' is a sentence of TFL. Although we can identify sentences of English when we encounter them, we do not have a formal definition of `sentence of English'. But in this chapter, we will \emph{define} exactly what will count as a sentence of TFL. This is one respect in which a formal language like TFL is more precise than a natural language like English.


\section{Expressions}

We have seen that there are three kinds of symbols in TFL:
\begin{center}
\begin{tabular}{l l}
Atomic sentences & $A,B,C,\ldots,Z,\ldots$\\
with subscripts, as needed & $A_1, B_1,Z_1,A_2,A_{25},J_{375},\ldots$\\[12pt]
Connectives & $\enot,\eand,\eor,\eif,\eiff$\\[12pt]
Brackets &( , )
\end{tabular}
\end{center}
Define an \define{expression of TFL} as any string of symbols of TFL.
So: write down any sequence of symbols of TFL, in any order, and you
have an expression of TFL.\footnote{In \cref{sec:contradiction}, we
will introduce the symbol~`$\ered$', which will also count as an
atomic sentence.}

\section{Sentences}\label{s:Sentences}
Given what we just said, `$(A \eand B)$' is an expression of TFL, and so is `$\lnot)(\eor()\eand(\enot\enot())((B$'. But the former is a \emph{sentence}, and the latter is \emph{gibberish}. We want some rules to tell us which TFL expressions are sentences.

Obviously, individual sentence letters like `$A$' and `$G_{13}$' should count as sentences. (We'll also call them \define{atomic} sentences.) We can form further sentences out of these by using the various connectives. Using negation, we can get `$\enot A$' and `$\enot G_{13}$'. Using conjunction, we can get `$(A \eand G_{13})$', `$(G_{13} \eand A)$', `$(A \eand A)$', and `$(G_{13} \eand G_{13})$'. We could also apply negation repeatedly to get sentences like `$\enot \enot A$' or apply negation along with conjunction to get sentences like `$\enot(A \eand G_{13})$' and `$\enot(G_{13} \eand \enot G_{13})$'. The possible combinations are endless, even starting with just these two sentence letters, and there are infinitely many sentence letters. So there is no point in trying to list all the sentences one by one.

Instead, we will describe the process by which sentences can be \emph{constructed}. Consider negation: Given any sentence \metav{A} of TFL, $\enot\metav{A}$ is a sentence of TFL. (Why the funny fonts? We return to this in \cref{s:Metavariables}.)

We can say similar things for each of the other connectives. For instance, if \metav{A} and \metav{B} are sentences of TFL, then $(\metav{A}\eand\metav{B})$ is a sentence of TFL. Providing clauses like this for all of the connectives, we arrive at the following formal definition for a \define{sentence of TFL}:
	\factoidbox{\label{TFLsentences}
	\begin{factoidnumber}
		\item Every sentence letter is a sentence.
		\item If \metav{A} is a sentence, then $\enot\metav{A}$ is a sentence.
		\item If \metav{A} and \metav{B} are sentences, then $(\metav{A}\eand\metav{B})$ is a sentence.
		\item If \metav{A} and \metav{B} are sentences, then $(\metav{A}\eor\metav{B})$ is a sentence.
		\item If \metav{A} and \metav{B} are sentences, then $(\metav{A}\eif\metav{B})$ is a sentence.
		\item If \metav{A} and \metav{B} are sentences, then $(\metav{A}\eiff\metav{B})$ is a sentence.
		\item Nothing else is a sentence.
	\end{factoidnumber}
	}
\newglossaryentry{sentence of TFL}
{
name=sentence (of TFL),
text = sentence of TFL,
plural = sentences of TFL,
description={A string of symbols in TFL that can be built up according to the inductive rules given \ifHTMLtarget in \cref{s:Sentences}\else on p.~\pageref{TFLsentences}\fi}
}

Definitions like this are called \define{inductive}. Inductive definitions begin with some specifiable base elements, and then present ways to generate indefinitely many more elements by compounding together previously established ones. To give you a better idea of what an inductive definition is, we can give an inductive definition of the idea of \emph{an ancestor of mine}. We specify a base clause:
	\begin{itemize}
		\item My parents are ancestors of mine.
	\end{itemize}
and then offer further clauses like:
	\begin{itemize}
		\item If $x$ is an ancestor of mine, then $x$'s parents are ancestors of mine.
	\end{itemize}
Finally, we stipulate that \emph{only} what the base and inductive
clauses say are ancestors of mine will count as such.
	\begin{itemize}
		\item Nothing else is an ancestor of mine.
	\end{itemize}
Using this definition, we can easily check to see whether someone is my ancestor: just check whether she is the parent of the parent of\ldots one of my parents. And the same is true for our inductive definition of sentences of TFL. Just as the inductive definition allows complex sentences to be built up from simpler parts, the definition allows us to decompose sentences into their simpler parts. Once we get down to sentence letters, then we know we are ok.

Let's consider some examples.

Suppose we want to know whether or not `$\enot \enot \enot D$' is a sentence of TFL. Looking at the second clause of the definition, we know that `$\enot \enot \enot D$' is a sentence \emph{if} `$\enot \enot D$' is a sentence. So now we need to ask whether or not `$\enot \enot D$' is a sentence. Again looking at the second clause of the definition, `$\enot \enot D$' is a sentence \emph{if} `$\enot D$' is. So, `$\enot D$' is a sentence \emph{if} `$D$' is a sentence. Now `$D$' is a sentence letter of TFL, so we know that `$D$' is a sentence by the first clause of the definition. So for a compound sentence like `$\enot \enot \enot D$', we must apply the definition repeatedly. Eventually we arrive at the sentence letters from which the sentence is built up.

Next, consider the example `$\enot (P \eand \enot (\enot Q \eor R))$'. Looking at the second clause of the definition, this is a sentence if `$(P \eand \enot (\enot Q \eor R))$' is, and this is a sentence if \emph{both} `$P$' \emph{and} `$\enot (\enot Q \eor R)$' are sentences. The former is a sentence letter, and the latter is a sentence if `$(\enot Q \eor R)$' is a sentence. It is. Looking at the fourth clause of the definition, this is a sentence if both `$\enot Q$' and `$R$' are sentences, and both are!

Ultimately, every sentence is constructed nicely out of sentence letters. When we are dealing with a \emph{sentence} other than a sentence letter, we can see that there must be some sentential connective that was introduced \emph{last}, when constructing the sentence. We call that connective the \define{main logical operator} of the sentence. In the case of `$\enot\enot\enot D$', the main logical operator is the very first `$\enot$' sign. In the case of `$(P \eand \enot (\enot Q \eor R))$', the main logical operator is `$\eand$'. In the case of `$((\enot E \eor F) \eif \enot\enot G)$', the main logical operator is `$\eif$'.

As a general rule, you can find the main logical operator for a sentence by using the following method:
\begin{itemize}
	\item If the first symbol in the sentence is `$\enot$', then that is the main logical operator
	\item Otherwise, start counting the brackets. For each open-bracket, i.e., `(', add $1$; for each closing-bracket, i.e., `$)$', subtract $1$. When your count is at exactly $1$, the first operator you hit (\emph{apart} from a `$\enot$') is the main logical operator.
\end{itemize}

(Note: if you do use this method, then make sure to include \emph{all} the brackets in the sentence, rather than omitting some as per the conventions of \cref{TFLconventions}!)

The inductive structure of sentences in TFL will be important when we consider the circumstances under which a particular sentence would be true or false. The sentence `$\enot \enot \enot D$' is true if and only if the sentence `$\enot \enot D$' is false, and so on through the structure of the sentence, until we arrive at the atomic components. We will return to this point in \cref{ch.TruthTables}.

The inductive structure of sentences in TFL also allows us to give a formal definition of the \emph{scope} of a negation (mentioned in \cref{s:ConnectiveConjunction}). The scope of a `$\enot$' is the subsentence for which `$\enot$' is the main logical operator. Consider a sentence like:
$$(P \eand (\enot (R \eand B) \eiff Q))$$
which was constructed by conjoining `$P$' with `$ (\enot (R \eand B) \eiff Q)$'. This last sentence was constructed by placing a biconditional between `$\enot (R \eand B)$' and `$Q$'. The former of these sentences---a subsentence of our original sentence---is a sentence for which `$\enot$' is the main logical operator. So the scope of the negation is just `$\enot(R \eand B)$'. More generally:
	\factoidbox{The \define{scope} of a connective (in a sentence) is the subsentence for which that connective is the main logical operator.}


\section{Bracketing conventions}
\label{TFLconventions}
Strictly speaking, the brackets in `$(Q \eand R)$' are an indispensable part of the sentence. Part of this is because we might use `$(Q \eand R)$' as a subsentence in a more complicated sentence. For example, we might want to negate `$(Q \eand R)$', obtaining `$\enot(Q \eand R)$'. If we just had `$Q \eand R$' without the brackets and put a negation in front of it, we would have `$\enot Q \eand R$'. It is most natural to read this as meaning the same thing as `$(\enot Q \eand R)$', but as we saw in  \cref{s:ConnectiveConjunction}, this is very different from `$\enot(Q\eand R)$'.

Strictly speaking, then, `$Q \eand R$' is \emph{not} a sentence. It is a mere \emph{expression}.

When working with TFL, however, it will make our lives easier if we are sometimes a little less than strict. So, here are some convenient conventions.

First,  we allow ourselves to omit the \emph{outermost} brackets of a sentence. Thus we allow ourselves to write `$Q \eand R$' instead of the sentence `$(Q \eand R)$'. However, we must remember to put the brackets back in, when we want to embed the sentence into a more complicated sentence!

Second, it can be a bit painful to stare at long sentences with many nested pairs of brackets. To make things a bit easier on the eyes, we will allow ourselves to use square brackets, `[' and~`]', instead of rounded ones. So there is no logical difference between `$(P\eor Q)$' and `$[P\eor Q]$', for example.

Combining these two conventions, we can rewrite the unwieldy sentence
$$(((H \eif I) \eor (I \eif H)) \eand (J \eor K))$$
rather more clearly as follows:
$$\bigl[(H \eif I) \eor (I \eif H)\bigr] \eand (J \eor K)$$
The scope of each connective is now much easier to pick out.

\practiceproblems

\solutions
\problempart
\label{pr.wiffTFL}
For each of the following: (a) Is it a sentence of TFL, strictly speaking? (b) Is it a sentence of TFL, allowing for our relaxed bracketing conventions?
\begin{compactlist}
\item $(A)$
\item $J_{374} \eor \enot J_{374}$
\item $\enot \enot \enot \enot F$
\item $\enot \eand S$
\item $(G \eand \enot G)$
\item $(A \eif (A \eand \enot F)) \eor (D \eiff E)$
\item $[(Z \eiff S) \eif W] \eand [J \eor X]$
\item $(F \eiff \enot D \eif J) \eor (C \eand D)$
\end{compactlist}

\problempart
Are there any sentences of TFL that contain no sentence letters? Explain your answer.\\

\problempart
What is the scope of each connective in the sentence
$$\bigl[(H \eif I) \eor (I \eif H)\bigr] \eand (J \eor K)$$

\chapter{Ambiguity}\label{s:AmbiguityTFL}

In English, sentences can be \define{ambiguous}, i.e., they can have more than one meaning.  There are many sources of ambiguity. One is \emph{lexical ambiguity:} a sentence can contain words which have more than one meaning.  For instance, `bank' can mean the bank of a river, or a financial institution. So I might say that `I went to the bank' when I took a stroll along the river, or when I went to deposit a check.  Depending on the situation, a different meaning of `bank' is intended, and so the sentence, when uttered in these different contexts, expresses different meanings.

A different kind of ambiguity is \emph{structural ambiguity}.  This arises when a sentence can be interpreted in different ways, and depending on the interpretation, a different meaning is selected.  A famous example due to Noam Chomsky is the following:
\begin{enumerate}
	\item[] Flying planes can be dangerous.
\end{enumerate}
There is one reading in which `flying' is used as an adjective which modifies `planes'. In this sense, what's claimed to be dangerous are airplanes which are in the process of flying.  In another reading, `flying' is a gerund: what's claimed to be dangerous is the act of flying a plane.  In the first case, you might use the sentence to warn someone who's about to launch a hot air baloon.  In the second case, you might use it to counsel someone against becoming a pilot.

When the sentence is uttered, usually only one meaning is intended. Which of the possible meanings an utterance of a sentence intends is determined by context, or sometimes by how it is uttered (which parts of the sentence are stressed, for instance). Often one interpretation is much more likely to be intended, and in that case it will even be difficult to ``see'' the unintended reading.  This is often the reason why a joke works, as in this example from Groucho Marx:
\begin{compactlist}
	\item[] One morning I shot an elephant in my pajamas.
	\item[] How he got in my pajamas, I don't know.
\end{compactlist}

Ambiguity is related to, but not the same as, vagueness. An adjective, for instance `rich' or `tall,' is \define{vague} when it is not always possible to determine if it applies or not.  For instance, a person who's 6~ft 4~in (1.9~m) tall is pretty clearly tall, but a building that size is tiny.  Here, context has a role to play in determining what the clear cases and clear non-cases are (`tall for a person,' `tall for a basketball player,' `tall for a building'). Even when the context is clear, however, there will still be cases that fall in a middle range.

In TFL, we generally aim to avoid ambiguity. We will try to give our symbolization keys in such a way that they do not use ambiguous words or  disambiguate them if a word has different meanings. So, e.g., your symbolization key will need two different sentence letters for `Rebecca went to the (money) bank' and `Rebecca went to the (river) bank.' Vagueness is harder to avoid. Since we have stipulated that every case (and later, every valuation) must make every basic sentence (or sentence letter) either true or false and nothing in between, we cannot accommodate borderline cases in TFL.

It is an important feature of sentences of TFL that they \emph{cannot} be structurally ambiguous. Every sentence of TFL can be read in one, and only one, way. This feature of TFL is also a strength. If an English sentence is ambiguous, TFL can help us make clear what the different meanings are.  Although we are pretty good at dealing with ambiguity in everyday conversation, avoiding it can sometimes be terribly important. Logic can then be usefully applied: it helps philosophers express their thoughts clearly, mathematicians to state their theorems rigorously, and software engineers to specify loop conditions, database queries, or verification criteria unambiguously.

Stating things without ambiguity is of crucial importance in the law as well. Here, ambiguity can, without exaggeration, be a matter of life and death. Here is a famous example of where a death sentence hinged on the interpretation of an ambiguity in the law. Roger Casement (1864--1916) was a British diplomat who was famous in his time for publicizing human-rights violations in the Congo and Peru (for which he was knighted in 1911). He was also an Irish nationalist. In 1914--16, Casement secretly travelled to Germany, with which Britain was at war at the time, and tried to recruit Irish prisoners of war to fight against Britain and for Irish independence. Upon his return to Ireland, he was captured by the British and tried for high treason.

The law under which Casement was tried is the \textit{Treason Act of 1351}. That act specifies what counts as treason, and so the prosecution had to establish at trial that Casement's actions met the criteria set forth in the Treason Act. The relevant passage stipulated that someone is guilty of treason
\begin{quote}
	if a man is adherent to the King's enemies in his
realm, giving to them aid and comfort in the realm, or elsewhere.
\end{quote}
Casement's defense hinged on the last comma in this sentence, which is not present in the original French text of the law from 1351.  It was not under dispute that Casement had been `adherent to the King's enemies', but the question was whether being adherent to the King's enemies constituted treason only when it was done in the realm, or also when it was done abroad. The defense argued that the law was ambiguous. The claimed ambiguity hinged on whether `or elsewhere' attaches only to `giving aid and comfort to the King's enemies' (the natural reading without the comma), or to both `being adherent to the King's enemies' and `giving aid and comfort to the King's enemies' (the natural reading with the comma).  Although the former interpretation might seem far fetched, the argument in its favor was actually not unpersuasive. Nevertheless, the court decided that the passage should be read with the comma, so Casement's antics in Germany were treasonous, and he was sentenced to death. Casement himself wrote that he was `hanged by a comma'.

We can use TFL to symbolize both readings of the passage, and thus to provide a disambiguiation. First, we need a symbolization key:
\begin{ekey}
	\item[A] Casement was adherent to the King's enemies in the realm
	\item[G] Casement gave aid and comfort to the King's enemies in the realm
	\item[B] Casement was adherent to the King's enemies abroad
	\item[H] Casement gave aid and comfort to the King's enemies abroad
\end{ekey}
The interpretation according to which Casement's behavior was not treasonous is this:
\[A \lor (G \lor H)\]
The interpretation which got him executed, on the other hand, can be symbolized by:
\[(A \lor B) \lor (G \lor H)\]
Remember that in the case we're dealing with, Casement was adherent to the King's enemies abroad ($B$~is true), but not in the realm, and he did not give the King's enemies aid or comfort in or outside the realm ($A$, $G$, and~$H$ are false).

One common source of structural ambiguity in English arises from its lack of parentheses. For instance, if I say `I like movies that are not long and boring', you will most likely think that what I dislike are movies that are long and boring. A less likely, but possible, interpretation is that I like movies that are both (a) not long and (b) boring. The first reading is more likely because who likes boring movies? But what about `I like dishes that are not sweet and flavorful'? Here, the more likely interpretation is that I like savory, flavorful dishes.  (Of course, I could have said that better, e.g., `I like dishes that are not sweet, yet flavorful'.) Similar ambiguities result from the interaction of `and' with `or'. For instance, suppose I ask you to send me a picture of a small and dangerous or stealthy animal.  Would a leopard count? It's stealthy, but not small. So it depends whether I'm looking for small animals that are dangerous or stealthy (leopard doesn't count), or whether I'm after either a small, dangerous animal or a stealthy animal (of any size).

These kinds of ambiguities are called \emph{scope ambiguities}, since they depend on whether or not a connective is in the scope of another. For instance, the sentence, `\textit{Avengers: Endgame} is not long and boring' is ambiguous between:
\begin{enumerate}
	\item\label{scamb1} \textit{Avengers: Endgame} is not: both long and boring.
	\item\label{scamb2} \textit{Avengers: Endgame} is both: not long and boring.
\end{enumerate}
\Cref*{scamb2} is certainly false, since \textit{Avengers: Endgame} is over three hours long. Whether you think~\cref*{scamb1} is true depends on if you think it is boring or not. We can use the symbolization key:
\begin{ekey}
	\item[B] \textit{Avengers: Endgame} is boring
	\item[L] \textit{Avengers: Endgame} is long
\end{ekey}
\Cref*{scamb1} can now be symbolized as `$\enot(L \eand B)$', whereas \cref*{scamb2} would be `$\enot L \eand B$'. In the first case, the `\eand' is in the scope of `\enot', in the second case `\enot' is in the scope of `\eand'.

The sentence `Tai Lung is small and dangerous or stealthy' is ambiguous between:
\begin{enumerate}
	\item\label{scamb3} Tai Lung is either both small and dangerous or stealthy.
	\item\label{scamb4} Tai Lung is both small and either dangerous or stealthy.
\end{enumerate}
We can use the following symbolization key:
\begin{ekey}
	\item[D] Tai Lung is dangerous
	\item[S] Tai Lung is small
	\item[T] Tai Lung is stealthy
\end{ekey}
The symbolization of \cref*{scamb3} is `$(S \eand D) \eor T$' and that of \cref*{scamb4} is `$S \eand (D \eor T)$'. In the first, `\eand' is in the scope of `\eor', and in the second `\eor' is in the scope of `\eand'.

\practiceproblems
\solutions
\problempart The following sentences are ambiguous. Give symbolization keys for each and symbolize the different readings.
\begin{compactlist}
	\item Haskell is a birder and enjoys watching cranes.
	\item The zoo has lions or tigers and bears.
	\item The flower is not red or fragrant.
\end{compactlist}

\chapter{Use and mention}\label{s:UseMention}
In this Part, we have talked a lot \emph{about} sentences. So we should pause to explain an important, and very general, point.

\section{Quotation conventions}
Consider these two sentences:
	\begin{itemize}
		\item Aristotle is the founding father of logic.
		\item The expression `Aristotle' is composed of one uppercase letter and eight lowercase letters.
	\end{itemize}
When we want to talk about the founding father of logic, we \emph{use}
his name. When we want to talk about the name of the founding father
of logic, we \emph{mention} that name, which we do by putting it in
quotation marks.

There is a general point here. When we want to talk about things in the world, we just \emph{use} words. When we want to talk about words, we typically have to \emph{mention} those words. We need to indicate that we are mentioning them, rather than using them. To do this, some convention is needed. We can put them in quotation marks, or display them centrally in the page (say). So this sentence:
	\begin{itemize}
		\item `Aristotle' was an ancient Greek philosopher.
	\end{itemize}
says that some \emph{expression} was an ancient Greek philosopher.
That's false. The \emph{man} was a Greek philosopher; his \emph{name} wasn't. Conversely,
this sentence:
	\begin{itemize}
		\item Aristotle is composed of one uppercase letter and eight lowercase letters.
	\end{itemize}
also says something false: Aristotle is a man (or was when he was alive), made of flesh rather than letters. One final example:
	\begin{itemize}
		\item ``\,`Aristotle'\,'' is the name of `Aristotle'.
	\end{itemize} 
On the left-hand-side, here, we have the name of a name. On the right hand side, we have a name. Perhaps this kind of sentence only occurs in logic textbooks, but it is true nonetheless.

Those are just general rules for quotation, and you should observe them carefully in all your work! To be clear, the quotation-marks here do not indicate reported speech. They indicate that you are moving from talking about an object, to talking about a name of that object.


\section{Object language and metalanguage}
These general quotation conventions are very important for us. After all, we are describing a formal language here, TFL, and so we must often \emph{mention} expressions from TFL.

When we talk about a language, the language that we are talking about is called the \define{object language}. The language that we use to talk about the object language is called the \define{metalanguage}.
\label{def.metalanguage}
\newglossaryentry{object language}
{
name=object language,
description={A language that is constructed and studied by logicians. In this textbook,
 the object languages are TFL and FOL}
}

\newglossaryentry{metalanguage}
{
name=metalanguage,
description={The language logicians use to talk about the object language. In this textbook, the metalanguage is English, supplemented by certain symbols like metavariables and technical terms like ``valid''}
}

For the most part, the object language in this chapter has been the formal language that we have been developing: TFL. The metalanguage is English. Not conversational English exactly, but English supplemented with some additional vocabulary to help us get along.

Now, we have used uppercase letters as sentence letters of TFL:
	$$A, B, C, Z, A_1, B_4, A_{25}, J_{375},\ldots$$
These are sentences of the object language (TFL). They are not sentences of English. So we must not say, for example:
	\begin{itemize}
		\item $D$ is a sentence letter of TFL.
	\end{itemize}
Obviously, we are trying to come out with an English sentence that says something about the object language (TFL), but `$D$' is a sentence of TFL, and not part of English. So the preceding is gibberish, just like:
	\begin{itemize}
		\item \foreignlanguage{german}{Schnee ist wei\ss} is a German sentence.
	\end{itemize}
What we surely meant to say, in this case, is:
	\begin{itemize}
		\item `\foreignlanguage{german}{Schnee ist wei\ss}' is a German sentence.
	\end{itemize}
Equally, what we meant to say above is just:
	\begin{itemize}
		\item `$D$' is a sentence letter of TFL.
	\end{itemize}
The general point is that, whenever we want to talk in English about some specific expression of TFL, we need to indicate that we are \emph{mentioning} the expression, rather than \emph{using} it. We can either deploy quotation marks, or we can adopt some similar convention, such as  placing it centrally in the page.


\section{Metavariables}\label{s:Metavariables}
However, we do not just want to talk about \emph{specific} expressions of TFL. We also want to be able to talk about \emph{any arbitrary} sentence of TFL. Indeed, we had to do this in \cref{s:Sentences}, when we presented the inductive definition of a sentence of TFL. We used uppercase script letters to do this, namely:
	$$\metav{A}, \metav{B}, \metav{C}, \metav{D}, \ldots$$
These symbols do not belong to TFL. Rather, they are part of our (augmented) metalanguage that we use to talk about \emph{any} expression of TFL. To explain why we need them, recall the second clause of the recursive definition of a sentence of TFL:
	\begin{compactlist}
		\item[2.] If $\metav{A}$ is a sentence, then $\enot \metav{A}$ is a sentence.
	\end{compactlist}
This talks about \emph{arbitrary} sentences. If we had instead offered:
	\begin{compactlist}
		\item[2$'$.] If `$A$' is a sentence, then `$\enot A$' is a sentence.
	\end{compactlist}
this would not have allowed us to determine whether `$\enot B$' is a sentence. To emphasize:
	\factoidbox{
	  `$\metav{A}$' is a symbol (called a \define{metavariable}) in augmented English, which we use to talk about expressions of TFL. 	`$A$' is a particular sentence letter of TFL.}

        \newglossaryentry{metavariables}
{
name=metavariables,
description={A variable in the metalanguage that can represent any sentence in the object language}
}
But this last example raises a further complication, concerning quotation conventions. We did not include any quotation marks in the second clause of our inductive definition. Should we have done so?

The problem is that the expression on the right-hand-side of this rule, i.e., `$\enot\metav{A}$', is not a sentence of English, since it contains~`$\enot$'. So we might try to write:
	\begin{numberlist}
		\item[2$''$.] If \metav{A} is a sentence, then `$\enot \metav{A}$' is a sentence.
	\end{numberlist}
But this is no good: `$\enot \metav{A}$' is not a TFL sentence, since `$\metav{A}$' is a symbol of (augmented) English rather than a symbol of TFL.

What we really want to say is something like this:
	\begin{numberlist}
		\item[2$'''$.] If \metav{A} is a sentence, then the result of concatenating the symbol `$\enot$' with the sentence \metav{A} is a sentence.
	\end{numberlist}
This is impeccable, but rather long-winded. %Quine introduced a convention that speeds things up here. In place of (2$''$), he suggested:
%	\begin{numberlist}
%		\item[2$'''$.] If \metav{A} and \metav{B} are sentences, then $\ulcorner (\metav{A}\eand\metav{B})\urcorner$ is a sentence
%	\end{numberlist}
%The rectangular quote-marks are sometimes called `Quine quotes', after Quine. The general interpretation of an expression like `$\ulcorner (\metav{A}\eand\metav{B})\urcorner$' is in terms of rules for concatenation.
But we can avoid long-windedness by creating our own conventions. We can perfectly well stipulate that an expression like `$\enot \metav{A}$' should simply be read \emph{directly} in terms of rules for concatenation. So, \emph{officially}, the metalanguage expression `$\enot \metav{A}$'
simply abbreviates:
\begin{quote}
	the result of concatenating the symbol `$\enot$' with the sentence \metav{A}
\end{quote}
and similarly, for expressions like `$(\metav{A} \eand \metav{B})$', `$(\metav{A} \eor \metav{B})$', etc.


\section{Quotation conventions for arguments}
One of our main purposes for using TFL is to study arguments, and that will be our concern in \cref{ch.TruthTables}. In English, the premises of an argument are often expressed by individual sentences, and the conclusion by a further sentence. Since we can symbolize English sentences, we can symbolize English arguments using TFL.

Or rather, we can use TFL to symbolize each of the \emph{sentences} used in an English argument. However, TFL itself has no way to flag some of them as the \emph{premises} and another as the \emph{conclusion} of an argument.  (Contrast this with natural English, which uses words like `so', `therefore', etc., to mark that a sentence is the \emph{conclusion} of an argument.)

%So, if we want to symbolize an \emph{argument} in TFL, what are we to do? 

%An obvious thought would be to add a new symbol to the \emph{object} language of TFL itself, which we could use to separate the premises from the conclusion of an argument. However, adding a new symbol to our object language would add significant complexity to that language, since that symbol would require an official syntax.\footnote{\emph{The following footnote should be read only after you have finished the entire book!} And it would require a semantics. Here, there are deep barriers concerning the semantics. First: an object-language symbol which adequately expressed `therefore' for TFL would not be truth-functional. (\emph{Exercise}: why?) Second: a paradox known as `validity Curry' shows that FOL itself \emph{cannot} be augmented with an adequate, object-language `therefore'.} 

So, we need another bit of notation. Suppose we want to symbolize the premises of an argument with $\metav{A}_1$, \dots,~$\metav{A}_n$ and the conclusion with $\metav{C}$. Then we will write:
$$\metav{A}_1, \ldots, \metav{A}_n \therefore \metav{C}$$
The role of the symbol `$\therefore$' is simply to indicate which sentences are the premises and which is the conclusion.

%Strictly, this extra notation is \emph{unnecessary}. After all, we could always just write things down long-hand, saying: the premises of the argument are well symbolized by $\metav{A}_1, \ldots \metav{A}_n$, and the conclusion of the argument is well symbolized by $\metav{C}$. But having some convention will save us some time. Equally, the particular convention we chose was fairly \emph{arbitrary}. After all, an equally good convention would have been to underline the conclusion of the argument. Still, this is the convention we will use.

Strictly, the symbol `$\therefore$' will not be a part of the object language, but of the \emph{metalanguage}. As such, one might think that we would need to put quote-marks around the TFL-sentences which flank it. That is a sensible thought, but adding these quote-marks would make things harder to read. Moreover---and as above---recall that \emph{we} are stipulating some new conventions. So, we can simply stipulate that these quote-marks are unnecessary. That is, we can simply write:
$$A, A \eif B \therefore B$$
\emph{without any quotation marks}, to indicate an argument whose premises are (symbolized by) `$A$' and `$A \eif B$' and whose conclusion is (symbolized by)~`$B$'.
